% !TeX root = ../main.tex
\chapter{Interfacce, Lambdas e Inner Classes}

\section{Interfacce in Java}

Un'interfaccia è un costrutto della programmazione ad oggetti che definisce un \textbf{contratto} di comportamento. Mentre una classe descrive \textit{cosa un oggetto è} (stato e comportamento), un'interfaccia descrive \textit{cosa un oggetto sa fare}.
\vspace{\baselineskip}

Un'interfaccia viene dichiarata con la parola chiave \texttt{interface}. Per sua natura, è una struttura puramente astratta:

\begin{itemize}
    \item \textbf{Metodi}: Tutti i metodi dichiarati in un'interfaccia sono implicitamente \texttt{public} e \texttt{abstract} (fino a Java 7). Non è necessario specificare questi modificatori.
    \item \textbf{Campi}: Qualsiasi variabile dichiarata in un'interfaccia è implicitamente \texttt{public}, \texttt{static} e \texttt{final} (costante).
    \item \textbf{Istanziazione}: Non è possibile creare un'istanza di un'interfaccia tramite \texttt{new}. Tuttavia, un'interfaccia può essere utilizzata come \textbf{tipo di dato} per riferimenti a oggetti di classi che la implementano. In tale contesto, il polimorfismo ed il pattern matching funzionano in modo del tutto analogo a quanto visto nel capitolo precedente.
\end{itemize}
\vspace{\baselineskip}

Una classe "aderisce" al contratto dell'interfaccia usando la parola chiave \texttt{implements}. 

\begin{itemize}
    \item La classe deve fornire un'implementazione concreta (\textit{override}) per \textbf{tutti} i metodi astratti dell'interfaccia, a meno che la classe stessa non sia dichiarata \texttt{abstract}. (N.B.: essendo i metodi di un'interfaccia implicitamente \texttt{public}, essi devono essere implementati con la stessa visibilità o superiore: non possono essere \texttt{protected} o \texttt{private}).
    \item \textbf{Ereditarietà Multipla}: A differenza delle classi, Java permette a una classe di implementare un numero illimitato di interfacce, risolvendo il problema dell'ereditarietà multipla.
\end{itemize}
\vspace{\baselineskip}

Oltre alle proprietà standard appena citate, le interfacce moderne presentano caratteristiche avanzate che ne estendono la potenza:

\begin{itemize}
    \item \textbf{Metodi static o private} (Java 8+): Un'interfaccia può fornire simili metodi di utilità (ovviamente non astratti), ma non possono essere invocati su istanze di classi che implementano l'interfaccia. Devono essere chiamati direttamente sull'interfaccia stessa.
    \item \textbf{Estensione tra interfacce}: Un'interfaccia può estendere una o più altre interfacce utilizzando la parola chiave \texttt{extends}. A differenza delle classi, qui l'ereditarietà multipla è permessa: un'interfaccia "figlia" eredita tutti i contratti delle "madri".
    \item \textbf{Implementazione in Records ed Enum}: Anche i \textbf{records} (introdotti stabilmente in Java 16) e le \textbf{enum} possono implementare interfacce. Questo permette, ad esempio, di definire un comportamento comune per un insieme di costanti o per oggetti che sono puri contenitori di dati.
    \item \textbf{Sealed Interfaces} (Java 17+): È possibile dichiarare un'interfaccia come \texttt{sealed}. In questo caso, lo sviluppatore può specificare esattamente quali classi o interfacce sono autorizzate a implementarla tramite la clausola \texttt{permits}, garantendo un controllo totale sulla gerarchia.
\end{itemize}


\subsection{L'interfaccia \texttt{Comparable}}
L'interfaccia \texttt{Comparable} è uno degli esempi più chiari di come Java utilizzi le interfacce per definire protocolli di interazione tra oggetti.

\begin{lstlisting}[language=Java]
public interface Comparable {
    int compareTo(Object other);
}
\end{lstlisting}

Il "contratto" di questa interfaccia si sviluppa su tre livelli distinti:

\begin{itemize}
    \item \textbf{Obbligo Sintattico}: A livello puramente formale, il compilatore richiede solo che la classe concreta implementi il metodo \texttt{compareTo} con la firma corretta e che restituisca un valore di tipo \texttt{int}. Il compilatore \textbf{non entra nel merito} della logica interna: se il metodo restituisse sempre \texttt{42}, il codice sarebbe sintatticamente corretto e verrebbe compilato senza errori, pur essendo logicamente inutile.
    \item \textbf{Obbligo Semantico}:  Molti algoritmi (come \texttt{Arrays.sort}) sono programmati assumendo che lo sviluppatore abbia rispettato la logica del confronto.
    Se la logica di \texttt{compareTo} è incoerente (es. $x < y$ ma il metodo dice il contrario), gli algoritmi di ordinamento produrranno risultati errati, loop infiniti o eccezioni a runtime. Si parla in tal caso di \textbf{contratto semantico}: l'interfaccia garantisce la forma, lo sviluppatore garantisce il significato.
    \item \textbf{La Logica del Valore di Ritorno}: La convenzione sui valori di ritorno è basata sul segno numerico, il che permette una gestione matematica molto efficiente del confronto:

        \begin{itemize}
            \item \textbf{Valore Negativo ($< 0$)}: Indica che \texttt{x} (l'oggetto corrente) è più piccolo di \texttt{y} (\texttt{other}). L'oggetto \texttt{x} deve apparire \textit{prima} di \texttt{y}.
            \item \textbf{Zero ($= 0$)}: Indica che gli oggetti sono equivalenti ai fini dell'ordinamento.
            \item \textbf{Valore Positivo ($> 0$)}: Indica che \texttt{x} è più grande di \texttt{y}. L'oggetto \texttt{x} deve apparire \textit{dopo} di \texttt{y}.
        \end{itemize}

\end{itemize}

\begin{tcolorbox}[colback=blue!5!white,colframe=blue!75!black,title=Nota: Il Contratto Semantico]
Il \textbf{contratto semantico} rappresenta l'insieme di vincoli e assunzioni sul comportamento di un modulo che non possono essere espressi tramite la sintassi del linguaggio (il "contratto sintattico").

In Java, molte librerie sono scritte assumendo che lo sviluppatore conosca e rispetti queste convenzioni. Un esempio classico è il legame tra \texttt{equals()} e \texttt{hashCode()}:
\begin{itemize}
    \item Se due oggetti sono uguali secondo \texttt{equals()}, \textbf{devono} restituire lo stesso \texttt{hashCode()}.
\end{itemize}
La violazione di questo contratto non genera errori di compilazione, ma porta a comportamenti imprevedibili e bug logici difficili da individuare (es. oggetti "scomparsi" da una \texttt{HashSet}).

Per quanto riguarda le interfacce, in Java sono lo strumento principale per definire il contratto sintattico, ma la documentazione Javadoc è il luogo dove viene sancito il contratto semantico.
\end{tcolorbox}



\subsubsection*{Esempio Pratico: La Classe \texttt{Employee}}
Supponiamo di voler rendere una classe "ordinabile" tramite i metodi di sistema (come \texttt{Arrays.sort}); essa deve innanzitutto implementare l'interfaccia \texttt{Comparable}: si utilizza la parola chiave \texttt{implements} nella firma della classe per comunicare al compilatore che la classe si impegna a rispettare il contratto di \texttt{Comparable}:

\begin{lstlisting}[language=Java]
class Employee implements Comparable {
    // ...
}
\end{lstlisting}

La classe deve poi fornire il corpo del metodo \texttt{compareTo}. Supponendo di voler ordinare i dipendenti in base al salario (\texttt{salary}), l'implementazione sarà la seguente:

\begin{lstlisting}[language=Java]
public int compareTo(Object otherObject) {
   // Cast obbligatorio: l'interfaccia accetta Object, 
   // noi dobbiamo trattarlo come Employee
   Employee other = (Employee) otherObject;
   
   // Utilizzo del metodo statico di utilita'
   return Double.compare(this.salary, other.salary);
}
\end{lstlisting}

Da notare come, invece di sottrarre manualmente i salari (operazione rischiosa con i numeri a virgola mobile a causa della precisione), abbiamo usato il metodo statico \texttt{Double.compare(x, y)}. Esso restituisce automaticamente $-1, 0, 1$ rispettando perfettamente il contratto semantico dell'interfaccia.

Grazie a questa implementazione, se avessimo un array di dipendenti chiamato \texttt{staff}, potremmo semplicemente scrivere \texttt{Arrays.sort(staff)}. L'algoritmo di ordinamento chiamerà internamente il metodo \texttt{compareTo} su ogni coppia di dipendenti per decidere come spostarli.


\subsubsection*{Evoluzione: l'interfaccia \texttt{Comparable} Generica}
Sebbene i generics saranno trattati in un capitolo successivo, vale la pena accennarne l'uso nel contesto specifico che stiamo trattando.

Nell'esempio precedente, l'interfaccia \texttt{Comparable} è stata utilizzata in una forma "grezza" (\textit{raw type}), il che comporta alcuni svantaggi significativi in termini di sicurezza dei tipi e leggibilità del codice; nello specifico, poiché il metodo dell'interfaccia accetta un \texttt{Object}, abbiamo dovuto eseguire un cast esplicito a \texttt{Employee} per poter accedere al campo \texttt{salary}. Se avessimo passato un oggetto di tipo diverso (es. una \texttt{String}), il codice avrebbe lanciato una \texttt{ClassCastException}.

A partire da Java 5, l'interfaccia \texttt{Comparable} è stata trasformata in un \textbf{tipo generico}. Questo aggiornamento permette di specificare immediatamente con quale tipo di oggetti la classe è in grado di confrontarsi:

\begin{lstlisting}[language=Java]
public interface Comparable<T> {
   int compareTo(T other);
}
\end{lstlisting}

Possiamo quindi migliorare la nostra classe \texttt{Employee} implementando la versione generica di \texttt{Comparable}, evitando così il fastidioso downcasting:

\begin{lstlisting}[language=Java]
class Employee implements Comparable<Employee>
{
   public int compareTo(Employee other)
   {
      return Double.compare(salary, other.salary);
   }
   . . .
}
\end{lstlisting}

Il vantaggio principale di questa evoluzione è che il compilatore ora può garantire la coerenza dei tipi durante la compilazione, eliminando la possibilità di errori di tipo a runtime e migliorando la leggibilità del codice.

\begin{important}
    L'uso di \texttt{Comparable<T>} è oggi lo standard assoluto. La versione non generica (\texttt{raw type}) viene mantenuta solo per compatibilità con il codice scritto prima del 2004, ma il suo utilizzo è fortemente sconsigliato (\textit{deprecated} nella logica di progettazione).
\end{important}

\subsubsection*{Comparable ed Ereditarietà: Il problema dell'Antisimmetria}
Il contratto di \texttt{Comparable} impone una regola matematica ferrea definita \textbf{antisimmetria}:
\begin{equation}
    sgn(x.compareTo(y)) = -sgn(y.compareTo(x))
\end{equation}
In termini semplici: se invertiamo gli addendi del confronto, il segno del risultato deve invertirsi. Se $x < y$, allora deve essere vero che $y > x$.
\vspace{\baselineskip}

I problemi insorgono quando una sottoclasse (es. \texttt{Manager}) estende una superclasse che implementa \texttt{Comparable} (es. \texttt{Employee}). 

Se \texttt{Manager} esegue l'override di \texttt{compareTo} tentando di confrontare campi specifici dei manager, si rischia di rompere l'antisimmetria:
\begin{itemize}
    \item \texttt{impiegato.compareTo(manager)}: Funziona, perché il manager viene trattato come un semplice impiegato (upcasting).
    \item \texttt{manager.compareTo(impiegato)}: Lancia una \texttt{ClassCastException} se il manager tenta di forzare l'impiegato a essere un manager per leggerne i campi extra.
\end{itemize}
\vspace{\baselineskip}

Esistono due strategie per gestire il confronto in una gerarchia di classi, a seconda del requisito di business:

\paragraph{Scenario A: Confronto vietato tra classi diverse}
Se classi diverse non devono mai essere confrontate tra loro (anche se una è sottoclasse dell'altra), il metodo \texttt{compareTo} deve iniziare con un controllo rigoroso sulla classe di appartenenza:
\begin{lstlisting}[language=Java]
if (getClass() != other.getClass()) 
    throw new ClassCastException();
\end{lstlisting}
Questo garantisce che il confronto avvenga solo tra "pari", evitando asimmetrie nel lancio di eccezioni.

\paragraph{Scenario B: Algoritmo di confronto comune}
Se esiste un modo logico per confrontare oggetti di diverse sottoclassi (es. una gerarchia aziendale), si deve adottare una strategia centralizzata:
\begin{enumerate}
    \item Definire il metodo \texttt{compareTo} come \textbf{final} nella superclasse.
    \item Utilizzare un metodo di supporto (es. \texttt{rank()}) che ogni sottoclasse può sovrascrivere per definire il proprio "peso" o posizione nell'ordinamento.
\end{enumerate}

\begin{example}
In una gerarchia \texttt{Employee}, \texttt{Manager}, \texttt{Executive}, si può definire un metodo \texttt{rank()} che restituisce un intero (es. 1 per segretari, 2 per manager, 3 per executive). Il \texttt{compareTo} nella superclasse userà questo valore per stabilire chi "viene prima", garantendo che la regola dei segni sia sempre rispettata.
\end{example}

\subsection{Interfacce vs Classi Astratte}

Una domanda sorge spontanea: perché i progettisti di Java hanno introdotto le interfacce se esistevano già le classi astratte? Non si poteva definire \texttt{Comparable} come una classe astratta?

\begin{lstlisting}[language=Java, caption={Perché questo approccio fallisce}]
// Se Comparable fosse una classe astratta...
abstract class Comparable {
   public abstract int compareTo(Object other);
}

// ...Employee non potrebbe estenderla se fosse gia figlio di Person
class Employee extends Person, Comparable // ERRORE DI COMPILAZIONE
\end{lstlisting}

Il motivo fondamentale è che Java permette a una classe di estendere **una sola** superclasse (ereditarietà singola). Se \texttt{Employee} estende già \texttt{Person}, non può estendere anche \texttt{Comparable}.



Tuttavia, una classe può implementare un numero arbitrario di interfacce:
\begin{lstlisting}[language=Java]
class Employee extends Person implements Comparable, Serializable, Cloneable // OK
\end{lstlisting}


I progettisti di Java (ispirandosi a C++) hanno scelto di non supportare l'ereditarietà multipla tra classi perché rende il linguaggio:
\begin{itemize}
    \item \textbf{Complesso}: si pensi al "diamante della morte" in C++ dove non si sa quale metodo della superclasse invocare.
    \item \textbf{Meno efficiente}: la gestione della memoria e dei puntatori diventa più onerosa.
\end{itemize}
Le interfacce offrono quasi tutti i benefici dell'ereditarietà multipla evitando però queste complicazioni, poiché non trasportano \textbf{stato} (variabili di istanza) ma solo \textbf{comportamento}.

\subsubsection*{Guida alla scelta: Classe Astratta o Interfaccia?}

La scelta tra una classe astratta e un'interfaccia non è solo tecnica, ma concettuale. Possiamo riassumere i criteri di scelta in questo modo:

\begin{itemize}
    \item \textbf{1. Identità vs. Comportamento (Is-a vs. Can-do)}:
    \begin{itemize}
        \item \textbf{Classe Astratta (Is-a)}: Si usa quando si vuole definire l'identità core di un oggetto. Una classe astratta risponde alla domanda: \textit{"Che cos'è questo oggetto?"}. Esempio: \texttt{Aereo} è un \texttt{Veicolo}.
        \item \textbf{Interfaccia (Can-do / Role)}: Si usa per definire un'abilità o un ruolo che un oggetto può ricoprire. Risponde alla domanda: \textit{"Cosa sa fare questo oggetto?"}. Esempio: \texttt{Aereo} è \texttt{Volante}, ma lo è anche un \texttt{Uccello} o un \texttt{Drone}.
    \end{itemize}
    \item \textbf{2. Stato vs. Puro Contratto}:
    \begin{itemize}
        \item Se hai bisogno di condividere dello \textbf{stato} (variabili di istanza non costanti) o costruttori comuni tra le sottoclassi, la \textbf{classe astratta} è l'unica scelta.
        \item Se vuoi definire solo un \textbf{contratto} comportamentale senza preoccuparti di come i dati verranno memorizzati, l' \textbf{interfaccia} è la scelta più pulita.
    \end{itemize}
\end{itemize}

\subsection{Metodi Default: Evoluzione e Compatibilità}
 Apartire da Java 8, utilizzando il modificatore \texttt{default} è possibile definire un metodo di istanza all'interno di un'interfaccia che include un'implementazione concreta: le classi che implementano l'interfaccia ereditano automaticamente il metodo senza l'obbligo di sovrascriverlo.

 Ciò permette l'evoluzione delle interfacce (aggiunta di nuovi metodi) senza rompere il codice esistente che le implementa (\textit{Binary Compatibility}). Tuttavia, ciò pone un problema: se una classe implementa due interfacce con metodi \texttt{default} identici (stessa firma), si verifica un conflitto di nomi. Java impone le seguenti regole per risolvere questo scenario:
\begin{enumerate}
    \item \textbf{Le classi vincono sulle interfacce}: Se una superclasse fornisce un'implementazione, il metodo \texttt{default} dell'interfaccia viene ignorato.
    \item \textbf{Override obbligatorio}: Se non c'è una superclasse a risolvere il conflitto, la classe deve obbligatoriamente implementare il metodo, decidendo quale comportamento adottare o fornendone uno nuovo.
\end{enumerate}

\section{Il Pattern Callback: Passare Comportamento}
Introduciamo ora un concetto che ci permetterà di comprendere appieno le Lambda Expressions: il pattern \textbf{callback}.

Nello studio della programmazione abbiamo imparato a passare dati (primitivi o oggetti) come parametri. Tuttavia, spesso nasce l'esigenza di passare non un dato, ma un \textbf{comportamento} (una funzione). In linguaggi come C o C++ è possibile passare direttamente un \textit{puntatore a funzione}. Si dice al programma: "Esegui il codice che si trova a questo indirizzo di memoria".
In Java, invece, le funzioni non esistono indipendentemente dalle classi. Per passare un'azione, dobbiamo "impacchettarla" dentro un riferimento a un oggetto. 

Le \textbf{Lambdas} (che vedremo a breve) sono semplicemente una "scorciatoia sintattica" per questo pattern. Senza aver compreso che in Java stiamo passando un \textit{riferimento a un oggetto che implementa un contratto}, la sintassi delle Lambda sembrerebbe magica. Capire la callback significa capire \textit{cosa succede dietro le quinte} di ogni Lambda Expression.


\textit{In sintesi: una callback è il modo in cui Java permette di trattare un blocco di codice come se fosse un oggetto, permettendone il passaggio tra diversi moduli del programma.}
\vspace{\baselineskip}

Il pattern \textbf{callback} (richiamata) si realizza come segue:
\begin{enumerate}
    \item Si definisce un'interfaccia con un singolo metodo (il "contratto").
    \item Si crea un oggetto che implementa quel metodo.
    \item Si passa il \textbf{riferimento} a tale oggetto.
\end{enumerate}
L'oggetto funge quindi da \textit{contenitore} (wrapper) per la funzione che vogliamo eseguire in differita.

\subsection{Un esempio di Callback: l'interfaccia \texttt{Comparator}}

Il pattern callback trova una sua tipica applicazione nella personalizzazione di algoritmi esistenti, come l'ordinamento.

\begin{itemize}
    \item \textbf{Il problema}: La classe \texttt{String} ha già un ordine naturale (alfabetico) definito tramite \texttt{Comparable}. Se vogliamo ordinare per lunghezza, non possiamo modificare la classe \texttt{String}.
    \item \textbf{La soluzione Callback}: Utilizziamo il metodo \texttt{Arrays.sort(T[] a, Comparator<? super T> c)}.
\end{itemize}

\begin{lstlisting}[language=Java, caption={Comparator come "pacchetto" di una funzione di confronto}]
// 1. Definiamo il comportamento (la funzione di callback)
class LengthComparator implements Comparator<String> {
   public int compare(String first, String second) {
      return first.length() - second.length();
   }
}

// 2. Passiamo l'oggetto/funzione all'algoritmo di sorting
Arrays.sort(friends, new LengthComparator());
\end{lstlisting}

Il metodo \texttt{Arrays.sort} è "cieco": sa come ordinare ma non sa \textit{cosa} rende un elemento maggiore di un altro. Passandogli un \texttt{Comparator}, gli stiamo fornendo la "vista" tramite una \textbf{callback}. Ogni volta che l'algoritmo deve confrontare due stringhe, "richiama" il metodo \texttt{compare} dell'oggetto che gli abbiamo fornito.



